\section{アプローチ12（複数フレーム情報を併用した時空間トラッキングおよび通行人誤検出防止手法）}

\subsection{手法の概要}
アプローチ11までの単一フレーム処理において判明した「待機列近傍を一時的に通過する通行人に対する誤検出・誤グループ化」を防ぐため、アプローチ12では前後の複数フレーム情報を組み合わせた動的トラッキング手法を導入した。前提として関心領域（ROI）処理により画面端の僅かな領域にしか通行人は映り込まない設計となっているが、本アプローチによりさらなる誤判定防止の頑健性を担保した。

処理システム全体におけるFPSを適応的に計測・維持するため、トップレベルに \texttt{measureFps.ts} を配置し、内部状態を共通保持させた。タイムスタンプの記録には、メモリ消費を最適化した64ビット浮動小数点数（\texttt{Float64Array}）によるリングバッファ構造（\texttt{ringBuffer.ts}）を使用している。

本アプローチにおける追跡判定および色空間を用いた特徴抽出の具体的仕様は以下の通りである。

\begin{itemize}
    \item \textbf{動的待機列判定とカルマンフィルタの適用}：
    人物検出の完了後かつグループ検出の前段階において、追跡配列を更新する。FPSから動的に割り出される所定の滞留時間閾値を超えて、一定割合（動的なバウンディングボックスサイズから算出される許容移動範囲）の位置に留まり続けている個体を「待機列構成員」として判定する。位置推定および補正にはカルマンフィルタ（\texttt{KalmanFilterForPeople.ts}）を導入した。
    
    \item \textbf{身体部位別カラー特徴抽出およびHSV比較}：
    追跡精度を強化するため、ポーズ検出で得られた身体部位（両肩・両腰で形成される胴体四角形の内側領域、および顔領域）からピクセル平均色を抽出・保持する。背景や手足など変化の激しい部分を除外し、カメラからの距離変化に対応できるよう領域切り出しサイズはバウンディングボックスの大きさに応じて動的に決定した。照明・影の影響を排除するため、色空間はRGBではなくHSVを使用し、彩度（S）の変動に連動して色相（H）と明度（V）の比較重みを動的に変更する判定ロジックを適用した。なお、切り出し画像の色抽出処理には \texttt{imageHelper} の固定サイズキャンバスを使い回すことで、メモリリークを完全に防止した。

    \item \textbf{TrackedPerson クラスによるライフサイクル管理}：
    各人物の追跡状態管理のため \texttt{TrackedPerson} クラスを追加した。コンストラクタで一意の ID（\texttt{crypto.randomUUID()}）を生成して \texttt{Person} 型に付与するとともに、FPSに連動した一定フレーム数以上見失った個体を使い回し可能な追跡配列から自動削除・解放する設計とした。
\end{itemize}

\subsection{結果およびアプローチ11との比較・課題}

\subsubsection{アプローチ11との比較}
時空間情報の統合、カルマンフィルタによる移動推測、およびHSV色判定の導入により、一時的に画面内を通過する通行人やカメラのブレ・ノイズによる瞬間的な誤検出が正確に排除された。これにより、実際に一定時間以上待機列内に滞留している人物のみを精緻に抽出してグループ化パイプラインへ引き渡すことが可能となった。

\subsubsection{判明した課題}
\begin{enumerate}
    \item \textbf{特徴量抽出・保持に伴う計算および実装コストの増大}：
    HSV色空間への変換、胴体四角形内領域の動的クロップ、および各部色の保持・更新といった色特徴量の計算処理が追加されたことで、全体としてのアルゴリズムおよびトラッキングロジックが極めて複雑化する課題が生じた。
    \item \textbf{特徴量非依存化の検証要求}：
    高度な移動判定（カルマンフィルタおよびFPS連動のバウンディングボックス移動割合チェック）のみで、重い特徴量抽出処理を廃止しても同等の追跡性能・通行人排除性能が維持できるかについての比較検証の必要性が浮彫りとなった。
\end{enumerate}
